\subsection{Composer 2: asynchronous RL for software-engineering agents
\citep{cursor2026composer2}}
\paragraph{Source type.} First-party technical report.
\paragraph{Research question.} Ask how continued pretraining, long-horizon environments,
asynchronous RL, and specialized infrastructure combine to train a practical coding agent.
\paragraph{Method.} Document full-parameter single-epoch policy-gradient training, the
modified group estimator, independent rollout and training workers, fast weight updates,
MoE routing replay, self-summarization, and the matching of training and deployment harnesses.
\paragraph{Main evidence.} Separate estimator and systems observations from the final
Composer benchmark results. Record evidence about best-of-$K$, long-horizon coherence,
training precision, load balancing, and inference--training consistency.
\paragraph{Connection to the position.} Make rollout staleness, trajectory compaction,
and realistic environment matching concrete parts of the post-training problem.
\paragraph{Critical assessment.} Examine the limited public recipe, private data and
benchmarks, coupled interventions, internal evaluation harness, and commercial source incentives.
\paragraph{Takeaway.} Complete after reading.
